\documentclass[11pt]{article}
\usepackage[margin=1.1in]{geometry}
\usepackage{amsmath,amssymb,amsthm,mathtools}
\usepackage{microtype}
\usepackage{lmodern}
\usepackage[T1]{fontenc}
\usepackage{xcolor}
\usepackage{booktabs}
\usepackage{enumitem}
\usepackage[colorlinks=true,linkcolor=blue!50!black,citecolor=blue!50!black,urlcolor=blue!50!black]{hyperref}

\newtheorem{theorem}{Theorem}
\newtheorem{proposition}[theorem]{Proposition}
\newtheorem{lemma}[theorem]{Lemma}
\newtheorem{corollary}[theorem]{Corollary}
\theoremstyle{remark}
\newtheorem{remark}[theorem]{Remark}

\newcommand{\N}{\mathbb{N}}
\newcommand{\Pp}{\mathbb{P}}
\newcommand{\E}{\mathbb{E}}
\DeclareMathOperator{\HH}{H}
\newcommand{\Pplus}{P^{+}}

\title{Runs of multiples in a finite set of integers:\\[2pt] $\sum k_i = \Theta(n \log n \log\log n)$}
\author{}
\date{7 October 2026}

\begin{document}
\maketitle

\begin{abstract}
Let $S$ be a set of $n$ positive integers.
For each $a\in S$, let $k(a)$ be the least positive integer $k$ such that $ka\notin S$.
We prove that $\max_{|S|=n}\sum_{a\in S}k(a) = \Theta(n\log n\log\log n)$.
The set $S=\{1,\dots,n\}$ gives only $\sim n\ln n$ and is therefore not extremal.
Sets of smooth numbers improve on it by a factor of order $\log\log n$.
For the upper bound, we apply an entropy argument to the prime factorisation.
The argument shows that primes below $(\log n)^4$ account for the entire excess over $n\log n$.
\end{abstract}

\section{Statement and proof overview}

Let $S$ be a finite set of positive integers, and write $|S|=n$.
For $a\in S$, the number $k(a)$ is the first multiplier for which $k(a)a$ is absent from $S$.
It will be convenient to work with the length $L(a)$ of the run of multiples that are present.
For each positive integer $t$, let $A_t$ be the set of elements whose runs have length at least $t$.
Thus
\[
  L(a) \coloneqq k(a)-1 = \max\{\ell : a,2a,\dots,\ell a\in S\} \ \ (\ge 1),
  \qquad
  A_t \coloneqq \{a\in S : L(a)\ge t\}.
\]
Counting the run lengths by their level sets gives
\[
  \sum_{a}k(a) = n+\sum_a L(a) = n + \sum_{t\ge1}|A_t|.
\]
Each run consists of distinct elements $a,2a,\dots,L(a)a$ of $S$, so $L(a)\le n$.

All logarithms are natural unless written $\log_2$.
We use $\Pplus(m)$ for the largest prime factor of $m$.
The letters $p,q$ always denote primes, and $\pi(x)$ counts the primes at most $x$.
We also write $\theta(x)=\sum_{p\le x}\log p$.

\begin{theorem}\label{thm:main}
There are absolute constants $0<c<C$ such that for all sufficiently large $n$,
\[
  c\, n\log n\log\log n \;\le\; \max_{|S|=n}\ \sum_{a\in S}k(a) \;\le\; C\, n\log n\log\log n .
\]
\end{theorem}

\paragraph{Proof overview.}
For the lower bound, we use a set of smooth numbers.
We choose its cutoff so that a fixed proportion of its elements begin long runs of multiples.
We then use disjoint scaled copies to obtain a set of every sufficiently large size.

For the upper bound, we expose the prime exponents in increasing order and apply the entropy chain rule.
A one-step entropy estimate first gives a bound on the number of pairs obtained by multiplying by a prime.
We then weight these pairs by the logarithm of the prime and split them into four classes.
Three classes can be bounded directly.
The fourth consists of primes in the top half of a run, so we avoid it by counting only primes in the bottom half.

\section{Lower bound: smooth numbers}

A positive integer is \emph{$y$-smooth} if all its prime factors are at most $y$.
For $y\ge 2$ and $x\ge1$, write $\Psi(x)=\Psi(x,y)=\#\{m\le x:\Pplus(m)\le y\}$ for the number of such integers up to $x$.
Our construction is the set
\[
  S = S(X,y) = \{m\le X : \Pplus(m)\le y\}.
\]

\paragraph{Why smooth numbers give long runs.}
Suppose that $a\in S$ and $a\le X/y$.
Every $j\le y$ satisfies both $\Pplus(j)\le y$ and $aj\le X$.
It follows that $aj\in S$, and hence $L(a)\ge y$.
There are $\Psi(X/y)$ choices of $a$ with this property, so
\begin{equation}\label{eq:lb-basic}
  \sum_{a\in S}L(a) \;\ge\; y\cdot \Psi(X/y).
\end{equation}
We therefore seek a cutoff $X$ for which $\Psi(X/y)\ge c\,\Psi(X)$ and $\log\Psi(X)\asymp y/\log y$.
The first condition gives long runs for a fixed proportion of the set.
The second gives the required scale $y\asymp\log n\log\log n$.

\subsection{Choosing the cutoff}

\begin{proposition}\label{prop:lb}
There is an absolute constant $c>0$ with the following property.
For every sufficiently large $y$, one can choose $X$ so that, with $S=S(X,y)$ and $n=|S(X,y)|$, we have
\[
  \log n\asymp y/\log y
  \qquad\text{and}\qquad
  \sum_{a\in S}k(a)\ge c\,n\log n\log\log n.
\]
\end{proposition}

\begin{proof}
Consider the cutoffs $X_i=e^{y}y^{i}$ for $0\le i\le K\coloneqq\lceil y/\log y\rceil$.
We first bound the count at the largest cutoff.
This will force at least one pair of consecutive cutoffs to have comparable counts.

\paragraph{Step 1: bound $\Psi(X_K)$ by Rankin's trick.}
Set $\sigma=1/\log y$.
Then
\[
  \Psi(x)\le\sum_{\Pplus(m)\le y}(x/m)^{\sigma}=x^{\sigma}\prod_{p\le y}\bigl(1-p^{-\sigma}\bigr)^{-1}.
\]
The definition of $K$ gives $\log X_K\le 2y+\log y$.
Thus $x^\sigma\le e\cdot e^{2y/\log y}$ whenever $x\le X_K$.

To bound the product, split the primes at $\sqrt y$.
For $\sqrt y<p\le y$, we have $p^{-\sigma}\le e^{-1/2}$, so each factor is at most $2.55$.
For $p\le\sqrt y$, apply $1-e^{-u}\ge u/2$ on $[0,1]$.
Each factor is then at most $2\log y/\log p\le 3\log y$.
Combining these bounds gives
\[
  \log\Psi(X_K)\le 1+\frac{2y}{\log y}+\pi(y)\log 2.55+\pi(\sqrt y)\log(3\log y)\le C_1\frac{y}{\log y}.
\]

\paragraph{Step 2: choose consecutive cutoffs by pigeonhole.}
The $K$ ratios $\Psi(X_i)/\Psi(X_{i-1})$ have product at most $e^{C_1y/\log y}$, since $\Psi(X_0)\ge1$.
There are $K\ge y/\log y$ such ratios, so at least one is at most $e^{C_1}$.
Choose an index $i$ with this property and set $X=X_i$.
Writing $n=\Psi(X)$, we obtain
\[
  \Psi(X/y)=\Psi(X_{i-1})\ge e^{-C_1}n.
\]
Equation~\eqref{eq:lb-basic} now gives $\sum_a L(a)\ge e^{-C_1}\,y\,n$.

\paragraph{Step 3: relate the size of the set to $y$.}
Every squarefree product of primes $\le y/2$ is below $4^{y/2}<e^{y}\le X$.
These products belong to $S$, giving $n\ge 2^{\pi(y/2)}$.
Step 1 also gives $n\le\Psi(X_K)\le e^{C_1y/\log y}$.
Together, the two bounds show that $\log n\asymp y/\log y$.
More specifically, $\log n\le C_1 y/\log y$ implies
\[
  y\ge C_1^{-1}\log n\,(\log\log n-\log C_1).
\]
Substituting this into the estimate from Step 2 yields
\[
  \sum_a k(a)\ge\sum_a L(a)\ge e^{-C_1}y\,n\ge c\,n\log n\log\log n.
\]
\end{proof}

\subsection{Obtaining every sufficiently large size}

We next pass from the sizes supplied by Proposition~\ref{prop:lb} to an arbitrary large $n$.
The proposition gives sets $S_y$ of size $n_y$ satisfying
\[
  \log 2\cdot\pi(y/2)\le\log n_y\le C_1y/\log y.
\]
Given large $n$, choose $y$ with $C_1y/\log y\approx\frac12\log n$.
Then $n_y\le\sqrt n$, while $\log n_y\ge\delta\log n$ for an absolute $\delta>0$.

Take $\lfloor n/n_y\rfloor$ scaled copies $P_jS_y$, where the $P_j$ are distinct primes with $P_j>\max S_y$.
Pad their union to size $n$ with further primes different from all $P_j$.
To check that this preserves each run, fix an element $P_js$ of one copy.
Every multiple of $P_js$ is divisible by $P_j$.
No element of another copy or of the padding is divisible by $P_j$.
Thus $L(P_js)=L_{S_y}(s)$ exactly, where the subscript denotes the set in which the run is measured.
The padding contributes only nonnegative terms to the sum.
At least $n/2$ elements lie in the copies, so the proposition gives $\sum k\ge c'n\log n\log\log n$ for every large $n$.

\subsection{Heuristic constant and numerical examples}

To estimate the constant heuristically, take $X=e^{y}$.
The saddle-point asymptotics for $\Psi(x,y)$ give $\Psi(X/t)\approx t^{-\alpha}\Psi(X)$, where the saddle-point parameter $\alpha$ satisfies $y^{-\alpha}\approx\tfrac12$.
They also give $\log n\approx(\log 4)\,y/\log y$.
These estimates suggest
\[
  \sum_a k(a)\approx\tfrac12 ny\approx\frac{1}{4\ln 2}\,n\ln n\ln\ln n\approx 0.36\,n\ln n\ln\ln n.
\]

The table reports $\sum_a L(a)/(n\ln n)$.
For the interval $S=\{1,\dots,n\}$, this ratio is $1+(2\gamma-1)/\ln n+o(1/\ln n)\approx 1.01$, where $\gamma$ is Euler's constant.

\begin{center}
\begin{tabular}{lrr}
\toprule
$S$ & $n$ (or $\ln n$) & $\sum L/(n\ln n)$ \\
\midrule
$\{1,\dots,n\}$ & $600{,}000$ & $1.012$ \\
$23$-smooth $\le e^{23}$ (exact) & $261{,}598$ & $1.129$ \\
$31$-smooth $\le e^{23.25}$ (exact) & $624{,}354$ & $1.167$ \\
$97$-smooth $\le e^{97}$ (log-binned count) & $\ln n\approx 37$ & $1.39$ \\
$197$-smooth $\le e^{197}$ (log-binned count) & $\ln n\approx 65$ & $1.54$ \\
$397$-smooth $\le e^{397}$ (log-binned count) & $\ln n\approx 113$ & $1.78$ \\
\bottomrule
\end{tabular}
\end{center}
The ratio grows roughly linearly in $\log y$, or equivalently like $\log\log n$.
An improvement is already visible for very small sets.
The set $\{1,2,3,4,6\}$ has $\sum L=11$, compared with $10$ for $\{1,\dots,5\}$.
Exhaustive search for $n\le 12$, and simulated annealing up to $n=50$, find smooth-type sets that beat $\{1,\dots,n\}$ whenever such sets exist.

\section{Upper bound}

\subsection{Entropy along the prime factorisation}

We begin by organising the prime factorisation in the order in which the entropy argument will use it.
Fix a prime $q$.
Every positive integer $m$ has a unique factorisation $m=c\,r$ in which all prime factors of $c$ are $<q$ and all prime factors of $r$ are $\ge q$.
Write $c=c_q(m)$ and call it the \emph{$q$-smooth part} of $m$.
We call the remaining factor $r$ \emph{$q$-rough}.

In this section, $X$ denotes a random variable uniformly distributed on $S$.
Let $\HH$ denote Shannon entropy in nats, so that $\HH(X)=\log n$.
For a prime $p$, write $v_p(X)$ for its exponent in $X$.
The exponent vector $(v_p(X))_p$ determines $X$.
Moreover, conditioning on $(v_p(X))_{p<q}$ is equivalent to conditioning on $c_q(X)$.
Apply the entropy chain rule with the primes in increasing order.
Only the finitely many primes dividing an element of $S$ contribute nonzero terms, and we obtain
\begin{equation}\label{eq:chain}
  \log n \;=\; \sum_{q}\HH\bigl(v_q(X)\,\big|\,c_q(X)\bigr).
\end{equation}

To estimate a term in this sum, fix a value $c$ of the smooth part.
Conditional on $c_q(X)=c$, we have $X=c\,r$, where $r$ is uniform on the \emph{fiber}
\[
  R_c\coloneqq\{r : r \text{ is $q$-rough and } cr\in S\},\qquad \Pp\bigl(c_q(X)=c\bigr)=|R_c|/n .
\]
The fiber $R_c$ depends on the fixed prime $q$.
The next lemma bounds the entropy of a prime exponent when multiplication by that prime keeps a given proportion of the set inside the set.

\begin{lemma}[one-step entropy]\label{lem:step}
Let $R$ be a finite set of positive integers, let $q$ be a prime, and suppose that $G\subseteq R$ satisfies $qG\subseteq R$.
Put $\gamma=|G|/|R|$.
If $r$ is uniform on $R$, then
\[
  \HH\bigl(v_q(r)\bigr)\;\ge\;\frac{\gamma}{2}\log\frac{2}{\gamma}.
\]
\end{lemma}

\begin{proof}
We first bound the mass outside the most likely exponent, then use that bound to estimate the entropy.

\paragraph{Step 1: bound the mass outside the most likely exponent.}
For each exponent $k$, define $\mu_k=\Pp(v_q(r)=k)$ and $b_k=\Pp(r\in G,\ v_q(r)=k)$.
Multiplication by $q$ maps $G\cap\{v_q=k\}$ injectively into $\{v_q=k+1\}$.
Consequently, $b_k\le\mu_{k+1}$.
We also have $b_k\le\mu_k$, since $G\subseteq R$.

Choose $j$ to maximise $\mu_j$, and set $\varepsilon=1-\mu_j$.
Use the bound by $\mu_k$ for all terms except $k=j$, and the bound by $\mu_{j+1}$ for that term.
This gives
\[
  \gamma=\sum_k b_k\le\sum_{k\ne j}\mu_k+\mu_{j+1}\le 2\varepsilon.
\]

\paragraph{Step 2: convert this mass bound into entropy.}
Let $h$ denote the binary entropy function.
There are two cases, according to the size of the largest probability $\mu_j$.
\begin{enumerate}[label=\emph{Case \arabic*.},leftmargin=*]
\item If $\mu_j\le\frac12$, then
\[
  \HH\bigl(v_q(r)\bigr)\ge\log(1/\mu_j)\ge\log 2\ge h(\gamma/2).
\]
\item Otherwise, $\gamma/2\le\varepsilon<\frac12$.
Retaining only whether the event $\{v_q=j\}$ occurs gives
\[
  \HH\bigl(v_q(r)\bigr)\ge h(\varepsilon)\ge h(\gamma/2),
\]
where the last inequality uses the fact that $h$ increases on $[0,\frac12]$.
\end{enumerate}
In either case, applying $h(x)\ge x\log(1/x)$ gives the stated bound.
\end{proof}

We will call a pair consisting of an integer in a set and a prime an \emph{edge} if multiplying the integer by the prime gives another element of the set.
The chain rule and the one-step lemma give the following bound on the number of edges.

\begin{lemma}[edge count]\label{lem:edges}
For every finite set $T$ of positive integers,
\[
  \#\{(t,p): p \text{ prime},\ t\in T,\ tp\in T\}\;\le\;\frac{2}{\log 2}\,|T|\log|T| .
\]
\end{lemma}

\begin{proof}
Take $X$ uniform on $T$ and apply \eqref{eq:chain}.
Fix a prime $p$ and a fiber $R_c$ with respect to $p$.
The set $G=\{r\in R_c: rp\in R_c\}$ satisfies $pG\subseteq R_c$.
Lemma~\ref{lem:step}, together with $\log(2/\gamma)\ge\log 2$, therefore gives
\[
  \HH(v_p\mid c)\ge\frac{\log 2}{2}\,|G|/|R_c|.
\]
Here $\HH(v_p\mid c)$ denotes the entropy of the exponent conditional on the smooth part being $c$.
Averaging over the fibers gives
\[
  \HH\bigl(v_p(X)\mid c_p(X)\bigr)\ge\frac{\log2}{2}\cdot\#\{t\in T:tp\in T\}/|T|.
\]
Finally, sum over $p$ and use \eqref{eq:chain}.
\end{proof}

\begin{remark}
Han's inequality, in its edge-isoperimetric form in $\mathbb{Z}^d$, gives the sharp constant $1/(2\log 2)$ in Lemma~\ref{lem:edges}.
The argument below only needs an absolute constant.
\end{remark}

We now apply the edge bound to the run lengths.
If $a\in A_q$, then $aq\in S$, so Lemma~\ref{lem:edges} gives
\begin{equation}\label{eq:edges}
  \sum_q|A_q|\;\le\;\frac{2}{\log 2}\,n\log n .
\end{equation}
Combining this with $\pi(L)\gg L/\log L$ and $L\le n$ already gives $\sum_a L(a)=O(n\log^2 n)$.
To reach the desired bound, we must replace one factor $\log n$ by $\log\log n$.

\subsection{Proof of the upper bound}

\paragraph{Weights and fibers.}
The next step is to retain the logarithmic weight of each prime.
Assign weight $\log q$ to every pair $(a,q)$ with $q\le L(a)$, and let
\[
  \Phi\coloneqq\sum_{a\in S}\theta\bigl(L(a)\bigr)=\sum_q|A_q|\log q .
\]
For a fixed $q$, write $a=c\,r$ with $c=c_q(a)$.
Within each fiber, define
\[
  G_c\coloneqq\{r\in R_c : cr\in A_q\},\qquad \gamma_c\coloneqq|G_c|/|R_c| .
\]
Thus $\gamma_c$ is the proportion of the fiber whose elements give runs of length at least $q$.
If $L(cr)\ge q$, then $c\,(rq)\in S$.
Also, $rq$ is $q$-rough, so $qG_c\subseteq R_c$.
Lemma~\ref{lem:step} therefore applies in every fiber.

Call a fiber \emph{sparse} when $\gamma_c\le q^{-1/2}$ and \emph{dense} otherwise.
Set $Q_0=(\log n)^4$.
We divide the pairs $(a,q)$ with $q\le L(a)$ into four classes, according to the size of $q$, the density of the fiber, and the length of the run.
For each of the first three classes, we bound the sum of the weights $\log q$.
The fourth class will disappear when we restrict to primes in the bottom half of each run.

\begin{enumerate}[label=\textbf{Class \arabic*.},leftmargin=*,itemsep=4pt]
\item \emph{Small primes: $q<Q_0$.}
Here we use the edge bound directly, bounding each weight by $\log Q_0$.
By \eqref{eq:edges}, the total weight is at most
\[
  \log Q_0\sum_q|A_q|\le\frac{8}{\log 2}\,n\log n\log\log n.
\]

\item \emph{Large primes in sparse fibers: $q\ge Q_0$ and $\gamma_c\le q^{-1/2}$.}
The one-step entropy estimate now supplies the logarithmic weight.
Sparsity gives $\log(2/\gamma_c)\ge\frac12\log q$.
Multiplying the bound from Lemma~\ref{lem:step} by the probability of the fiber, we get
\[
  \Pp(c_q(X)=c)\,\HH(v_q\mid c)\;\ge\;\frac{|R_c|}{n}\cdot\frac{\gamma_c}{4}\log q=\frac{|G_c|\log q}{4n}.
\]
Sum over sparse fibers and all $q$, and apply \eqref{eq:chain}.
The total weight of this class is at most $4\,n\log n$.

\item \emph{Large primes in dense fibers with long runs: $q\ge Q_0$, the fiber is dense, and $L(a)\ge 2q$.}
Long runs produce edges for every prime between $q$ and $2q$.
We use the edge bound inside the fiber to limit how many such runs there can be.

Let $E_c=\{r\in G_c : L(cr)\ge 2q\}$.
For each $r\in E_c$ and each prime $p\in(q,2q]$, the run condition gives $c\,(rp)\in S$.
The integer $rp$ is $q$-rough, so $rp\in R_c$.
The resulting pairs are distinct edges of $R_c$.
Apply Lemma~\ref{lem:edges} within $R_c$, bounding the logarithm of the fiber size by $\log n$.
Density gives $|R_c|<\sqrt q\,|G_c|$, and hence
\[
  |E_c|\,\bigl(\pi(2q)-\pi(q)\bigr)\le\frac{2}{\log 2}|R_c|\log n<\frac{2}{\log 2}\sqrt q\,|G_c|\log n .
\]
For large $q$, the prime number theorem gives $\pi(2q)-\pi(q)\ge q/(2\log q)$.
Any bound $\pi(2q)-\pi(q)\gg q/\log q$, such as one from Chebyshev's estimates, would suffice with a change in constants.
It follows that
\[
  |E_c|\le\frac{4}{\log 2}|G_c|\log n\,\log q/\sqrt q.
\]
To sum the weights, multiply this estimate by $\log q$ and sum over the dense fibers and the primes in this class.
The function $\log^2q/\sqrt q$ decreases for $q\ge e^4$, so we can bound it by its value at $Q_0$.
Using \eqref{eq:edges}, the total weight is therefore at most
\[
  \frac{4}{\log2}\,\log n\cdot\frac{\log^2Q_0}{\sqrt{Q_0}}\sum_q|A_q| = O\bigl(n(\log\log n)^2\bigr).
\]

\item \emph{Large primes in dense fibers with short runs: $q\ge Q_0$, the fiber is dense, and $L(a)<2q$.}
In this class, $L(a)/2<q\le L(a)$.
We will not need an estimate for its weight.
\end{enumerate}

\paragraph{Count only the bottom half of each run.}
Every pair with $q\le L(a)/2$ belongs to one of the first three classes.
Summing their bounds therefore gives
\[
  \sum_{a\in S}\theta\bigl(L(a)/2\bigr)\;\le\;(\text{Class }1)+(\text{Class }2)+(\text{Class }3)\;=\;O(n\log n\log\log n).
\]
Here each class on the right denotes its total weight.
Chebyshev's estimate states that $\theta(x)\ge c_0x$ for $x\ge2$, with an absolute $c_0>0$.
Thus, whenever $L(a)\ge4$, we have
\[
  L(a)\le\frac{2}{c_0}\theta\bigl(L(a)/2\bigr).
\]
We absorb the shorter runs and the initial $n$ in the relation between $k(a)$ and $L(a)$ into the term $4n$ below.
This gives
\[
  \sum_{a\in S}k(a)=n+\sum_{a\in S}L(a)\le 4n+\frac{2}{c_0}\,O(n\log n\log\log n)=O(n\log n\log\log n),
\]
which proves the upper bound and completes Theorem~\ref{thm:main}.
\qed

\section{Remarks}

\begin{enumerate}[leftmargin=*,itemsep=4pt]
\item \emph{Where the $\log\log n$ comes from.}
Classes~2 and~3 show that pairs $(a,q)$ with $(\log n)^4\le q\le L(a)/2$ have total weight only $O(n\log n)$, as in $\{1,\dots,n\}$.
Class~4 contains primes in the top half of the run and only affects the constant in the final bound.
The additional factor $\log\log n$ enters through Class~1, which contains primes below $(\log n)^4$.
This is also the range used by the smooth-number construction: its primes are at most $\approx\log n\log\log n$.

\item \emph{What entropy adds to the edge count.}
In the edge bound \eqref{eq:edges}, the contribution of a prime $q$, after division by $n$, is only $\Pp(L(X)\ge q)$.
For $\{1,\dots,n\}$, this is $\approx 1/q$, whereas $\HH(v_q(X)\mid c_q(X))\approx(\log q)/q$.
The extra factor $\log q$ records information about where the boundary for multiplication by $q$ lies.
The chain rule in increasing prime order captures this factor whenever the fibers are sparse.
In dense fibers, almost every run stops before $2q$, as the negligible contribution of Class~3 shows.

\item \emph{Constants.}
We have not optimised the constants.
Using $\theta(x)\sim x$, the argument gives roughly $\sum k\le 23\,n\ln n\ln\ln n$.
The coefficient is about $6$ if we use Han's constant in \eqref{eq:edges}.
The construction gives the heuristic value $\approx0.36\,n\ln n\ln\ln n$.
The correct leading constant remains open.
\end{enumerate}

\end{document}
